\begin{abstract}
A terminal with a low Earth orbit (LEO) satellite path and a terrestrial 5G path meets packet deadlines only if it knows how late each path may deliver in the next seconds. From \DataHours{} hours of 5\,Hz round-trip time (RTT) measurements of Starlink, OneWeb, and 5G links, we ask what of this is predictable. Past RTT changes hardly predict the next change, and point forecasts beat persistence mainly by expecting a return to the recent level. What is predictable is how far the RTT may move. Its spread differs by up to an order of magnitude between quiet and loaded periods, and its errors are skewed upward. A constant-width interval around a point forecast is therefore too wide when a link is quiet and too narrow when it is volatile. We show that interpolating ping streams onto a common time base leaks future samples into the inputs, which leaves a trained forecaster with a 0.2\,s error \LeakOptDSMAEhOne\% higher in deployment than in its evaluation. We forecast quantiles with an anchored recurrent encoder-decoder and calibrate them with conformalized quantile regression on held-out sessions. Against the same network with a scalar output, the 95\% interval score improves on all six links by \BSISWorstAbs\% to \BSISBestAbs\%, with session-bootstrap confidence intervals that exclude zero. A linear quantile model recovers most of this gain, which identifies the calibrated distributional output as its source. In an emulated deadline-aware duplication scheduler, the calibrated bound needs \SchedSaveSLDL\% less duplication on Starlink downlink with no more deadline misses.
\end{abstract}
\begin{IEEEkeywords}
Conformal prediction, LEO satellite networks, link-state forecasting, multipath transport, OneWeb, quantile regression, Starlink, 5G.
\end{IEEEkeywords}
